\notepage{67}{Orbital magnetic moment in quantum theory}
A particle of mass $m$ and charge $q$ in a magnetic field
$\mathbf B(\mathbf R)=\boldsymbol\nabla_R\times\mathbf A$ has, in Gaussian
units,
\begin{align*}
 H&=\frac{(\mathbf P-q\mathbf A/c)^2}{2m}\\
  &=\underbrace{\frac{\mathbf P^2}{2m}}_{\text{zeroth order}}
    -\underbrace{\frac{q}{2mc}(\mathbf P\cdot\mathbf A+
                                \mathbf A\cdot\mathbf P)}_{\text{first order in }\mathbf A/c}
    +\underbrace{\frac{q^2\mathbf A^2}{2mc^2}}_{\text{second order}}.
\end{align*}
Choose the Coulomb gauge, $\boldsymbol\nabla\cdot\mathbf A=0$. With
$\braket{\mathbf r}{\psi}=\psi(\mathbf r)$,
\begin{align*}
 \mel{\mathbf r}{\mathbf P\cdot\mathbf A}{\psi}
 &=-\ii\hbar\boldsymbol\nabla\cdot\bigl(\mathbf A(\mathbf r)\psi(\mathbf r)\bigr)\\
 &=-\ii\hbar\bigl[(\boldsymbol\nabla\cdot\mathbf A)\psi
                         +\mathbf A\cdot\boldsymbol\nabla\psi\bigr]\\
 &=\mathbf A(\mathbf r)\cdot(-\ii\hbar\boldsymbol\nabla\psi)
  =\mel{\mathbf r}{\mathbf A\cdot\mathbf P}{\psi}.
\end{align*}
Thus $\mathbf P\cdot\mathbf A=\mathbf A\cdot\mathbf P$ in this gauge.
For a uniform field $\mathbf B=B\hat{\mathbf e}_z$, choose
\[
 \mathbf A=\frac B2(-Y\hat{\mathbf e}_x+X\hat{\mathbf e}_y),
 \quad\boldsymbol\nabla\cdot\mathbf A=0,
 \quad\boldsymbol\nabla\times\mathbf A=B\hat{\mathbf e}_z.
\]
The term linear in the field becomes
\begin{align*}
 H_1&=-\frac q{mc}\mathbf A\cdot\mathbf P
     =-\frac q{mc}\frac B2(-YP_x+XP_y)\\
    &=-\frac q{2mc}\mathbf L\cdot\mathbf B
     =-\boldsymbol\mu\cdot\mathbf B,
 \qquad \boldsymbol\mu=\frac q{2mc}\mathbf L.
\end{align*}
In an $L_z$ eigenstate,
\[
 \mu_z=\frac q{2mc}L_z,
 \qquad\mu_z\text{ has eigenvalues }\frac{q\hbar}{2mc}(0,\pm1,\pm2,\ldots).
\]
The corresponding positive magneton scales are
\[
 \mu_B=\frac{e\hbar}{2m_ec},\qquad
 \mu_N=\frac{e\hbar}{2m_pc}.
\]
\editorial{The standard nuclear magneton uses the proton mass $m_p$.
Replacing it by a general nucleon mass $M$, as in the source's mnemonic,
defines the analogous scale $e\hbar/(2Mc)$ rather than the conventional
definition of $\mu_N$. Here $e>0$; an electron has charge $q=-e$.}
